\documentclass[a4paper,12pt]{article}
\usepackage[T2A]{fontenc}
\usepackage[utf8]{inputenc}
\usepackage[russian]{babel}
\usepackage{amsmath,amssymb,mathtools}
\renewcommand{\familydefault}{\sfdefault}
\usepackage{icomma}
\usepackage[a4paper,margin=18mm,headsep=7mm,footskip=10mm]{geometry}
\usepackage{microtype}
\usepackage{tikz}
\usetikzlibrary{calc,arrows.meta,positioning,shapes.geometric}
\usepackage{tkz-euclide}
\usepackage{pgfplots}
\pgfplotsset{compat=1.18}
\usepackage{tabularx,booktabs}
\usepackage{enumitem}
\usepackage{parskip}
\usepackage{fancyhdr}
\usepackage{setspace}
\usepackage{xcolor}

\definecolor{accent}{HTML}{1B6F78}
\definecolor{darkaccent}{HTML}{124B52}
\definecolor{textgray}{HTML}{30363A}
\definecolor{softgray}{HTML}{EEF2F3}
\definecolor{warmgray}{HTML}{6A7175}

\onehalfspacing
\setlength{\parindent}{0pt}
\setlength{\headheight}{24.1pt}
\setlist[itemize]{leftmargin=6mm,itemsep=2pt,topsep=3pt}
\setlist[enumerate]{leftmargin=7mm,itemsep=4pt,topsep=4pt}
\newcommand{\checkitem}{\(\square\)\hspace{0.45em}}

\pagestyle{fancy}
\fancyhf{}
\fancyhead[L]{\small\color{warmgray}Тригонометрические функции}
\fancyhead[R]{\small\color{warmgray}Ярослав Гаврилов}
\fancyfoot[C]{\small\color{warmgray}\thepage}
\renewcommand{\headrulewidth}{0.3pt}
\renewcommand{\headrule}{\hbox to\headwidth{\color{softgray}\leaders\hrule height \headrulewidth\hfill}}

\newcommand{\sectionrule}{\par\vspace{-0.3em}{\color{accent}\rule{\linewidth}{0.7pt}}\vspace{0.6em}}
\newcommand{\key}[1]{\textbf{\color{darkaccent}#1}}
\newcommand{\R}{\mathbb{R}}

\begin{document}

% ---------- TITLE ----------
\thispagestyle{empty}
\begin{center}
\vspace*{18mm}
{\Large\color{warmgray}Чек-лист}\par
\vspace{7mm}
{\Huge\bfseries\color{darkaccent}Тригонометрические функции}\par
\vspace{3mm}
{\Large периодичность, графики, неравенства и отбор корней}\par
\vspace{18mm}
{\large Ярослав Гаврилов}\par
\vspace{4mm}
{\normalsize 28 сентября 2026}\par
\vfill
{\small\color{warmgray}Лёвин Артём Александрович эксклюзивно для Ярослава Гаврилова}\par
\vspace{9mm}
\begin{tikzpicture}[remember picture,overlay]
  \draw[accent,line width=1.2pt,opacity=.65]
    ($(current page.south west)+(14mm,13mm)$)
    .. controls ($(current page.south west)+(58mm,25mm)$) and ($(current page.south east)+(-70mm,2mm)$) ..
    ($(current page.south east)+(-14mm,14mm)$);
\end{tikzpicture}
\end{center}
\clearpage

% ---------- 1 ----------
\section*{1. Периодичность: что нужно уметь}
\sectionrule

\checkitem \key{Понимать определение.} Число $T\ne0$ является периодом функции $f$, если вместе с каждым $x\in D(f)$ числа $x+T$ и $x-T$ также принадлежат $D(f)$ и
\[
 f(x+T)=f(x).
\]
Для доказательства достаточно проверить сдвиг на $+T$. Равенство $f(x-T)=f(x)$ следует из того же свойства после замены аргумента; на занятии обе подстановки использовались как симметричная самопроверка.

\checkitem \key{Знать базовые периоды.}
\begin{center}
\begin{tabular}{@{}lc@{}}
\toprule
\textbf{Функция} & \textbf{Наименьший положительный период}\\
\midrule
$\sin x,\ \cos x$ & $2\pi$\\
$\tg x,\ \ctg x$ & $\pi$\\
\bottomrule
\end{tabular}
\end{center}

\checkitem \key{Учитывать коэффициент перед $x$.} Если $k\ne0$, то
\[
T_{\sin(kx)}=T_{\cos(kx)}=\frac{2\pi}{|k|},
\qquad
T_{\tg(kx)}=T_{\ctg(kx)}=\frac{\pi}{|k|}.
\]

\subsection*{Пример 1. Доказать период $\boldsymbol{\frac{2\pi}{3}}$ для $\boldsymbol{f(x)=\sin 3x}$}
\[
\begin{aligned}
f\!\left(x+\frac{2\pi}{3}\right)
&=\sin\!\left(3x+2\pi\right)=\sin 3x=f(x),\\
f\!\left(x-\frac{2\pi}{3}\right)
&=\sin\!\left(3x-2\pi\right)=\sin 3x=f(x).
\end{aligned}
\]
Следовательно, $\frac{2\pi}{3}$ --- период функции. По формуле $\frac{2\pi}{|3|}$ это также её наименьший положительный период.

\subsection*{Пример 2. Почему $\boldsymbol{3\pi}$ не является периодом $\boldsymbol{y=\cos\frac{x}{2}}$}
\[
\cos\frac{x+3\pi}{2}
=\cos\left(\frac{x}{2}+\frac{3\pi}{2}\right)
=\sin\frac{x}{2}.
\]
Получилось другое выражение, значит равенство $f(x+3\pi)=f(x)$ выполняется не для всех $x$. Поэтому $3\pi$ периодом не является. Правильный наименьший положительный период:
\[
T=\frac{2\pi}{1/2}=4\pi.
\]

\subsection*{Быстрая самопроверка}
\begin{itemize}
\item Коэффициент $k$ стоит \textbf{внутри} аргумента: $\sin(kx)$, $\cos(kx)$, $\tg(kx)$.
\item Чем больше $|k|$, тем меньше период.
\item Для $\tg$ и $\ctg$ базовый период $\pi$, для $\sin$ и $\cos$ --- $2\pi$.
\item Примеры из занятия: $T\!\left(\sin\frac{4x}{5}\right)=\frac{5\pi}{2}$, \quad $T\!\left(\tg\frac{3x}{5}\right)=\frac{5\pi}{3}$.
\end{itemize}

% ---------- 2 ----------
\section*{2. Область определения и область значений}
\sectionrule

\checkitem \key{Синус и косинус ограничены:}
\[
D(\sin x)=D(\cos x)=\R,\qquad
-1\le \sin x\le1,\quad -1\le\cos x\le1.
\]
То есть область значений обеих функций равна $[-1;1]$.

\checkitem \key{Тангенс и котангенс по значениям не ограничены:}
\[
E(\tg x)=E(\ctg x)=\R.
\]
При этом
\[
D(\tg x)=\R\setminus\left\{\frac{\pi}{2}+\pi n\right\},\qquad
D(\ctg x)=\R\setminus\{\pi n\},\quad n\in\mathbb Z.
\]

\subsection*{Как найти область значений $\boldsymbol{y=a\sin x+b}$ или $\boldsymbol{y=a\cos x+b}$}
Поскольку $-1\le\sin x\le1$ и $-1\le\cos x\le1$, удобно помнить итог:
\[
\boxed{\ b-|a|\le y\le b+|a|\ }.
\]

\textbf{Пример.} Для $y=3\cos x-2$:
\[
-1\le\cos x\le1
\quad\Longrightarrow\quad
-3\le3\cos x\le3
\quad\Longrightarrow\quad
-5\le3\cos x-2\le1.
\]
Следовательно, $E(y)=[-5;1]$.

\textbf{Пример со знаком «минус».} Для $y=-2\sin x+1$ можно либо аккуратно умножить двойное неравенство на $-2$ с разворотом знаков, либо сразу применить формулу:
\[
1-|-2|\le y\le1+|-2|
\quad\Longrightarrow\quad
-1\le y\le3.
\]

\subsection*{Типичные ошибки}
\begin{itemize}[itemsep=0pt,topsep=2pt]
\item при умножении двойного неравенства на отрицательное число забывают развернуть оба знака; после преобразований границы записывают от меньшего значения к большему;
\item путают $D(f)$ и $E(f)$ или ошибочно ограничивают значения тангенса промежутком $[-1;1]$.
\end{itemize}

\clearpage

% ---------- 3 ----------
\section*{3. График $\boldsymbol{y=2\sin x+1}$}
\sectionrule

\checkitem \key{Понимать преобразования.}
\begin{itemize}
\item множитель $2$ \textbf{перед} $\sin x$ растягивает график по вертикали в $2$ раза;
\item $+1$ сдвигает весь график вверх на $1$;
\item амплитуда становится равной $2$;
\item область значений: $[-1;3]$;
\item период остаётся $2\pi$, потому что коэффициент перед $x$ не изменился.
\end{itemize}

\checkitem \key{Уметь построить по ключевым точкам.}
\begin{center}
\begin{tabular}{@{}cccccc@{}}
\toprule
$x$ & $-\pi$ & $-\frac{\pi}{2}$ & $0$ & $\frac{\pi}{2}$ & $\pi$\\
\midrule
$2\sin x+1$ & $1$ & $-1$ & $1$ & $3$ & $1$\\
\bottomrule
\end{tabular}
\end{center}

\begin{center}
\begin{tikzpicture}
\begin{axis}[
  width=0.92\linewidth,
  height=7.0cm,
  axis lines=middle,
  xmin=-3.45,xmax=3.45,
  ymin=-2.2,ymax=3.8,
  xtick={-3.14159,-1.5708,0,1.5708,3.14159},
  xticklabels={$-\pi$,$-\frac{\pi}{2}$,$0$,$\frac{\pi}{2}$,$\pi$},
  ytick={-2,-1,0,1,2,3},
  grid=both,
  minor tick num=1,
  clip=false,
  legend style={at={(0.02,0.98)},anchor=north west,draw=none,fill=none},
  samples=220,
  domain=-3.3:3.3
]
\addplot[accent,very thick] {2*sin(deg(x))+1};
\addlegendentry{$y=2\sin x+1$}
\addplot[warmgray,dashed,thick] {sin(deg(x))};
\addlegendentry{$y=\sin x$}
\addplot[only marks,mark=*,mark size=2pt,accent] coordinates {
(-3.14159,1) (-1.5708,-1) (0,1) (1.5708,3) (3.14159,1)
};
\end{axis}
\end{tikzpicture}
\end{center}

\subsection*{Мини-алгоритм построения}
\begin{enumerate}
\item Возьмите стандартные узловые аргументы $-\pi,-\frac\pi2,0,\frac\pi2,\pi$.
\item Вычислите значения функции.
\item Отметьте точки на координатной плоскости.
\item Соедините их плавной синусоидой и продолжите график периодически.
\item Проверьте амплитуду, среднюю линию $y=1$ и период $2\pi$.
\end{enumerate}

% ---------- 4 ----------
\section*{4. Тригонометрические неравенства на окружности}
\sectionrule

\checkitem \key{Помнить геометрический смысл косинуса.} На единичной окружности $\cos\alpha$ --- это абсцисса точки. Поэтому условие на косинус выделяет вертикальной границей нужные дуги.

\subsection*{Пример 1. $\boldsymbol{\cos\alpha\ge\frac12}$}
Граничные углы: $-\frac\pi3$ и $\frac\pi3$. Нужна правая дуга, где абсцисса не меньше $\frac12$:
\[
\boxed{\alpha\in\left[-\frac\pi3+2\pi n;\ \frac\pi3+2\pi n\right],\quad n\in\mathbb Z.}
\]

\subsection*{Пример 2. $\boldsymbol{\cos\alpha\le\frac{\sqrt3}{2}}$}
Граничные углы: $\frac\pi6$ и $\frac{11\pi}{6}$. Нужна большая дуга, расположенная слева от вертикали $x=\frac{\sqrt3}{2}$:
\[
\boxed{\alpha\in\left[\frac\pi6+2\pi n;\ \frac{11\pi}{6}+2\pi n\right],\quad n\in\mathbb Z.}
\]

\begin{center}
\begin{minipage}[t]{0.46\linewidth}
\centering
\begin{tikzpicture}[scale=0.82]
  \tkzDefPoint(0,0){Oa}
  \tkzDefPoint(2.15,0){Aa}
  \tkzDrawCircle[accent,line width=1pt](Oa,Aa)
  \draw[->,warmgray] (-2.45,0)--(2.55,0);
  \draw[->,warmgray] (0,-2.4)--(0,2.5);
  \draw[darkaccent,dashed] (1.075,-2.0)--(1.075,2.0);
  \draw[accent,line width=3pt,line cap=round] ({2.15*cos(-60)},{2.15*sin(-60)}) arc[start angle=-60,end angle=60,radius=2.15];
  \fill[accent] ({2.15*cos(-60)},{2.15*sin(-60)}) circle (2pt);
  \fill[accent] ({2.15*cos(60)},{2.15*sin(60)}) circle (2pt);
  \node[below right] at ({2.15*cos(-60)},{2.15*sin(-60)}) {$-\frac\pi3$};
  \node[above right] at ({2.15*cos(60)},{2.15*sin(60)}) {$\frac\pi3$};
\end{tikzpicture}

{\small $\cos\alpha\ge\frac12$: правая дуга}
\end{minipage}\hfill
\begin{minipage}[t]{0.46\linewidth}
\centering
\begin{tikzpicture}[scale=0.82]
  \tkzDefPoint(0,0){Ob}
  \tkzDefPoint(2.15,0){Ab}
  \tkzDrawCircle[accent,line width=1pt](Ob,Ab)
  \draw[->,warmgray] (-2.45,0)--(2.55,0);
  \draw[->,warmgray] (0,-2.4)--(0,2.5);
  \draw[darkaccent,dashed] ({2.15*sqrt(3)/2},-2.0)--({2.15*sqrt(3)/2},2.0);
  \draw[accent,line width=3pt,line cap=round]
    ({2.15*cos(30)},{2.15*sin(30)}) arc[start angle=30,end angle=330,radius=2.15];
  \fill[accent] ({2.15*cos(30)},{2.15*sin(30)}) circle (2pt);
  \fill[accent] ({2.15*cos(330)},{2.15*sin(330)}) circle (2pt);
  \node[above right] at ({2.15*cos(30)},{2.15*sin(30)}) {$\frac\pi6$};
  \node[below right] at ({2.15*cos(330)},{2.15*sin(330)}) {$\frac{11\pi}{6}$};
\end{tikzpicture}

{\small $\cos\alpha\le\frac{\sqrt3}{2}$: большая дуга}
\end{minipage}
\end{center}

\subsection*{Что обязательно проверить}
\begin{itemize}
\item знаки $\ge$ и $\le$ включают граничные точки; при $>$ и $<$ скобки будут круглыми;
\item интервал записывается от меньшего угла к большему;
\item для полного ответа добавляется $2\pi n$, $n\in\mathbb Z$;
\item при движении по выделенной дуге удобно ориентироваться против часовой стрелки.
\end{itemize}

\clearpage

% ---------- FLOWCHART ----------
\thispagestyle{empty}
\begin{center}
{\Large\bfseries\color{darkaccent}Алгоритм: неравенство с $\cos\alpha$ на окружности}

\vspace{2mm}
{\small Схема ниже относится к основному случаю $|a|\le1$. Если $|a|>1$, сначала сравните $a$ с диапазоном $-1\le\cos\alpha\le1$.}
\end{center}
\vspace{5mm}

\begin{center}
\begin{tikzpicture}[
  node distance=11mm,
  startstop/.style={ellipse,draw=accent,line width=1pt,minimum width=42mm,minimum height=10mm,align=center,fill=softgray},
  process/.style={rectangle,rounded corners=2mm,draw=accent,line width=1pt,text width=100mm,minimum height=12mm,align=center},
  decision/.style={diamond,aspect=2.4,draw=accent,line width=1pt,text width=38mm,align=center,inner sep=1.5mm},
  arrow/.style={-{Stealth[length=3mm]},line width=1pt,draw=darkaccent}
]
  \node[startstop] (s) {Начало};
  \node[process,below=of s] (p1) {Найдите граничные углы из уравнения $\cos\alpha=a$};
  \node[process,below=of p1] (p2) {Отметьте соответствующие точки на единичной окружности};
  \node[decision,below=12mm of p2] (d) {Какой знак?};
  \node[process,below left=15mm and 8mm of d,text width=57mm] (p3a) {$\cos\alpha\ge a$ или $>a$: выберите дугу с большими абсциссами};
  \node[process,below right=15mm and 8mm of d,text width=57mm] (p3b) {$\cos\alpha\le a$ или $<a$: выберите дугу с меньшими абсциссами};
  \node[process,below=25mm of d] (p4) {Запишите выбранную дугу от меньшего угла к большему и добавьте $2\pi n$, $n\in\mathbb Z$};
  \node[process,below=of p4] (p5) {Проверьте тип скобок: строгий знак --- точки исключаются; нестрогий --- включаются};
  \node[startstop,below=of p5] (e) {Ответ};

  \draw[arrow] (s)--(p1);
  \draw[arrow] (p1)--(p2);
  \draw[arrow] (p2)--(d);
  \draw[arrow] (d)-- node[above left]{\small $\ge,>$} (p3a);
  \draw[arrow] (d)-- node[above right]{\small $\le,<$} (p3b);
  \draw[arrow] (p3a.south) |- (p4.west);
  \draw[arrow] (p3b.south) |- (p4.east);
  \draw[arrow] (p4)--(p5);
  \draw[arrow] (p5)--(e);
\end{tikzpicture}
\end{center}
\clearpage

% ---------- 5 ----------
\section*{5. Отбор корней на заданном промежутке}
\sectionrule

После решения тригонометрического уравнения часто получается серия
\[
x=x_0+2\pi n,\qquad n\in\mathbb Z.
\]
Нужно оставить только те значения, которые попадают в заданный промежуток.

\begin{minipage}[t]{0.48\linewidth}
\subsection*{Способ 1. Двойное неравенство}
Для промежутка $[a;b]$ запишите
\[
a\le x_0+2\pi n\le b.
\]
Решите неравенство относительно целого $n$, затем подставьте найденные значения обратно.

\textbf{Пример.} Для
\[
x=\frac\pi6+2\pi n
\]
на $\left[-\frac\pi2;\pi\right]$:
\[
-\frac\pi2\le\frac\pi6+2\pi n\le\pi,
\]
\[
-\frac13\le n\le\frac{5}{12}.
\]
Единственное целое значение $n=0$, поэтому
\[
\boxed{x=\frac\pi6}.
\]
\end{minipage}\hfill
\begin{minipage}[t]{0.48\linewidth}
\subsection*{Способ 2. По окружности}
\begin{enumerate}[leftmargin=6mm,itemsep=2pt]
\item Отметьте дугу заданного промежутка.
\item Нанесите базовый угол $x_0$.
\item Найдите его копии с шагом $2\pi$ внутри дуги.
\item Запишите конкретные углы без параметра $n$.
\end{enumerate}

\textbf{Характерный случай.} Для той же серии на
\[
\left[-\frac{5\pi}{2};-\frac{3\pi}{2}\right]
\]
подходит $n=-1$:
\[
\boxed{x=\frac\pi6-2\pi=-\frac{11\pi}{6}}.
\]

\subsection*{Контроль перед ответом}
\begin{itemize}[leftmargin=5mm,itemsep=1.5pt]
\item \checkitem серия корней записана верно;
\item \checkitem границы и скобки учтены;
\item \checkitem используются только целые $n$;
\item \checkitem найденные углы лежат в промежутке.
\end{itemize}
\end{minipage}

% ---------- TRAINING ----------
\clearpage
\section*{6. Тренировочный блок --- Путь А}
\sectionrule

\textbf{10 заданий.} Выполняйте по порядку: сложность постепенно возрастает. Решения здесь не приводятся.

\begin{enumerate}[itemsep=4pt,topsep=4pt]
\item Найдите наименьший положительный период функции $y=\sin 5x$.
\item Найдите наименьший положительный период функции $y=\tg\frac{2x}{3}$.
\item Определите, является ли $T=2\pi$ периодом функции $y=\cos\frac{x}{2}$. Кратко обоснуйте ответ подстановкой $x+T$.
\item Найдите область значений функции $y=4\cos x-1$.
\item Найдите область значений функции $y=-3\sin x+2$.
\item Для $y=2\sin x-1$ найдите значения при $x=-\pi,-\frac\pi2,0,\frac\pi2,\pi$ и укажите область значений функции.
\item Решите неравенство $\displaystyle \cos x\ge\frac{\sqrt2}{2}$.
\item Решите неравенство $\displaystyle \cos x\le\frac12$.
\item Из серии $\displaystyle x=\frac\pi4+2\pi n$, $n\in\mathbb Z$, отберите корни на $\displaystyle \left[-\frac\pi2;\pi\right]$.
\item Из серии $\displaystyle x=\frac\pi3+2\pi n$, $n\in\mathbb Z$, отберите корни на $\displaystyle \left[-\frac{5\pi}{2};-\frac{3\pi}{2}\right]$.
\end{enumerate}

\vspace{5mm}
\textbf{Мини-проверка после выполнения:} период берётся из базовой функции; область значений учитывает $|a|$; в неравенствах проверены границы дуги; при отборе корней параметр $n$ обязательно целый.

% ---------- ANSWERS ----------
\clearpage
\section*{7. Ответы}
\sectionrule

\renewcommand{\arraystretch}{1.45}
\begin{tabularx}{\linewidth}{@{}cX@{}}
\toprule
\textbf{№} & \textbf{Ответ}\\
\midrule
1 & $\displaystyle \frac{2\pi}{5}$\\
2 & $\displaystyle \frac{3\pi}{2}$\\
3 & Нет. $\displaystyle \cos\frac{x+2\pi}{2}=\cos\left(\frac{x}{2}+\pi\right)=-\cos\frac{x}{2}\not\equiv\cos\frac{x}{2}$.\\
4 & $[-5;3]$\\
5 & $[-1;5]$\\
6 & $-1,-3,-1,1,-1$ соответственно; $E(y)=[-3;1]$.\\
7 & $\displaystyle x\in\left[-\frac\pi4+2\pi n;\frac\pi4+2\pi n\right],\ n\in\mathbb Z$.\\
8 & $\displaystyle x\in\left[\frac\pi3+2\pi n;\frac{5\pi}{3}+2\pi n\right],\ n\in\mathbb Z$.\\
9 & $\displaystyle x=\frac\pi4$.\\
10 & $\displaystyle x=-\frac{5\pi}{3}$.\\
\bottomrule
\end{tabularx}

\vfill
{\small\color{warmgray}Лёвин Артём Александрович эксклюзивно для Ярослава Гаврилова}\par

\end{document}
